\documentclass[conference]{IEEEtran}
\usepackage{graphicx}
\usepackage{amsmath}
\usepackage{hyperref}

%\title{Fuzzy Logic-Based Maximum Power Point Tracking for PV Battery Chargers}

\title{Fuzzy Logic-Based MPPT with Advanced Buck Converter Modeling for PV Battery Chargers}

\begin{document}

\author{\IEEEauthorblockN{Mehmet Gök, Kemal Ozanoglu, and Günhan Dündar}
\IEEEauthorblockA{\small{Department of Electrical and Electronics Engineering, Bogazici University}}
}

\maketitle

\begin{abstract}
This paper presents a Fuzzy Logic Controller (FLC) for Maximum Power Point Tracking (MPPT) in photovoltaic (PV) battery charging systems. The proposed system utilizes an advanced buck converter to efficiently transfer power from the PV panel to the battery. The buck converter is designed to operate at its maximum efficiency point, ensuring optimal power delivery. The FLC dynamically adjusts the duty cycle of the buck converter to track the maximum power point (MPP) under varying solar irradiation and load conditions. Simulation results demonstrate the effectiveness of the proposed approach, showing improved efficiency and dynamic response compared to conventional MPPT methods. The paper also includes a detailed analysis of the buck converter's performance, with plots of duty cycle versus output current and efficiency. The results highlight the benefits of using a fuzzy logic-based MPPT system for PV battery chargers.
\end{abstract}

\begin{IEEEkeywords}
Fuzzy Logic Control, Maximum Power Point Tracking (MPPT), Photovoltaic (PV) Systems, Buck Converter, Battery Charging, Energy Efficiency
\end{IEEEkeywords}

\section{Introduction}
The increasing demand for renewable energy sources has led to the widespread adoption of photovoltaic (PV) systems for power generation. One of the key challenges in PV systems is maximizing the energy harvested from the solar panels, which is achieved through Maximum Power Point Tracking (MPPT). Traditional MPPT methods, such as Perturb and Observe (P\&O) and Incremental Conductance (IC), have limitations in terms of dynamic response and efficiency under rapidly changing environmental conditions. Fuzzy Logic Control (FLC) has emerged as a promising alternative due to its ability to handle nonlinear systems and uncertainties.

In this paper, we propose a Fuzzy Logic-based MPPT system for PV battery chargers using an advanced buck converter. The buck converter is designed to operate at its maximum efficiency point, ensuring optimal power transfer from the PV panel to the battery. The FLC dynamically adjusts the duty cycle to track the MPP under varying solar irradiation and load conditions. The proposed system is validated through simulations, demonstrating improved efficiency and dynamic response compared to conventional MPPT methods.

\section{System Design}
\subsection{Solar Panel and Buck Converter}
The solar panel is model 6MN6A280 from Ablytek with non-ideal characteristics, such as varying irradiation and temperature. The buck converter steps down the voltage from the PV panel to the level required by the battery while maintaining high efficiency. The design of the buck converter is optimized to operate at its maximum efficiency point, ensuring minimal power loss during the conversion process.

\begin{figure}[ht]
\centering
\includegraphics[width=0.8\columnwidth]{fuzzysystem.png}
\caption{Schematic of the advanced buck converter used in the PV battery charging system.}
\label{fig:buck_design}
\end{figure}

\section{Buck Converter Performance}
The buck converter is a critical component of the PV battery charging system. It steps down the voltage from the PV panel to the level required by the battery while maintaining high efficiency. The design of the buck converter is optimized to operate at its maximum efficiency point, ensuring minimal power loss during the conversion process. To evaluate the performance of the buck converter, a DC source is connected to the output to emulate a voltage source. The output current (\(I_{out}\)), PV efficiency, and converter efficiency are measured as a function of the duty cycle.

\subsection{Output Current vs Duty Cycle}
The relationship between the output current (\(I_{out}\)) and the duty cycle is critical for understanding the behavior of the buck converter. As the duty cycle increases, the output current initially rises, reaching a maximum value at the optimal duty cycle. Beyond this point, further increases in the duty cycle lead to a decrease in output current due to non-ideal effects such as switching losses and inductor saturation.

\begin{figure}[ht]
\centering
\includegraphics[width=0.8\columnwidth]{Iout_vs_D.png}
\caption{Duty cycle versus output current for the buck converter. The peak output current occurs at the optimal duty cycle, indicating the maximum power point (MPP) of the PV panel.}
\label{fig:duty_current}
\end{figure}

\subsection{PV Efficiency vs Duty Cycle}
The efficiency of the PV system is directly influenced by the duty cycle of the buck converter. At low duty cycles, the efficiency is suboptimal due to insufficient power transfer. As the duty cycle increases, the efficiency improves, reaching a maximum at the optimal duty cycle. Beyond this point, efficiency decreases due to increased losses in the converter.

\begin{figure}[ht]
\centering
\includegraphics[width=0.8\columnwidth]{PVefficency_vs_D.png}
\caption{Duty cycle versus PV efficiency for the buck converter. The peak efficiency corresponds to the optimal duty cycle, ensuring maximum power transfer from the PV panel to the battery.}
\label{fig:duty_pv_efficiency}
\end{figure}

\subsection{Converter Efficiency vs Duty Cycle}
The efficiency of the buck converter itself is also a key performance metric. Converter efficiency is defined as the ratio of output power to input power. At low duty cycles, the converter efficiency is low due to high conduction losses. As the duty cycle increases, the efficiency improves, reaching a maximum at the optimal duty cycle. Beyond this point, switching losses and inductor core losses dominate, leading to a decrease in efficiency.

\begin{figure}[ht]
\centering
\includegraphics[width=0.8\columnwidth]{convertereffic_vs_D.png}
\caption{Duty cycle versus converter efficiency for the buck converter. The peak converter efficiency aligns with the optimal duty cycle, ensuring minimal power loss during the conversion process.}
\label{fig:duty_converter_efficiency}
\end{figure}
\vspace{\baselineskip}
Hence according to above results maximizing PV efficiency will be enough for maximizing output current which is critical for battery charging applications.

\section{Fuzzy Inference}
The control loop model is critical for ensuring the system operates at the maximum power point (MPP). The Fuzzy Logic Controller (FLC) is designed to dynamically adjust the duty cycle of the buck converter to track the MPP. The FLC uses two input variables: the change in power (\(\Delta P\)) and the change in voltage (\(\Delta V\)). These inputs are defined as follows:

\begin{equation}
\Delta P = P(t) - P(t-1)
\end{equation}

\begin{equation}
\Delta V = V(t) - V(t-1)
\end{equation}

where:
\begin{itemize}
    \item \(P(t)\) is the instantaneous power at time \(t\),
    \item \(P(t-1)\) is the power at the previous time step,
    \item \(V(t)\) is the instantaneous voltage at time \(t\),
    \item \(V(t-1)\) is the voltage at the previous time step.
\end{itemize}
\vspace{\baselineskip}

\subsection{Fuzzy Logic Controller Design}
The FLC is designed to handle the nonlinear relationship between \(\Delta P\) and \(\Delta V\). The input variables are fuzzified using membership functions, as shown in the figures below. The output is determined based on a 5x5 fuzzy rule table, which maps the fuzzified inputs to the appropriate control action. These rules ensure that the FLC adjusts the duty cycle to move the operating point toward the MPP.

\subsubsection{Membership Functions}
The membership functions for \(\Delta P\) and \(\Delta V\) are defined using three distinct shapes: triangular, trapezoidal, and Gaussian. These functions map the crisp input values to fuzzy sets, such as \textit{Negative Large (NL)}, \textit{Negative Small (NS)}, \textit{Zero (Z)}, \textit{Positive Small (PS)}, and \textit{Positive Large (PL)}. The output membership function is similarly defined to represent the control action.(Figure 5,6,7.) 

\begin{figure}[ht]
\centering
\includegraphics[width=0.8\columnwidth]{DPmember.png}
\caption{Membership functions for \(\Delta P\). The input is fuzzified into five linguistic variables: Negative Large (NL), Negative Small (NS), Zero (Z), Positive Small (PS), and Positive Large (PL).}
\label{fig:membership_deltaP}
\end{figure}

\begin{figure}[ht]
\centering
\includegraphics[width=0.8\columnwidth]{DVmember.png}
\caption{Membership functions for \(\Delta V\). The input is fuzzified into five linguistic variables: Negative Large (NL), Negative Small (NS), Zero (ZE), Positive Small (PS), and Positive Large (PL).}
\label{fig:membership_deltaV}
\end{figure}

\begin{figure}[ht]
\centering
\includegraphics[width=0.8\columnwidth]{Dmember.png}
\caption{Membership functions for the output. The control action is represented using five linguistic variables: Negative Large (NL), Negative Small (NS), Zero (ZE), Positive Small (PS), and Positive Large (PL).}
\label{fig:membership_output}
\end{figure}

\subsubsection{Fuzzy Rules}
The fuzzy rules are designed to determine the appropriate control action based on the fuzzified inputs. A 5x5 rule table is used to map the combinations of \(\Delta p\) and \(\Delta V\) to the output. The complete rule table is at table 1.

\begin{table}[ht]
\centering
\caption{5x5 Fuzzy rule table for the FLC.}
\label{tab:fuzzy_rules}
\begin{tabular}{|c|c|c|c|c|c|}
\hline
\(\Delta V \backslash \Delta P\) & NL & NS & ZE & PS & PL \\ \hline
NL & NL & NL & NS & PS & PL \\ \hline
NS & NL & NS & NS & PS & PS \\ \hline
ZE & Z & Z & Z & Z & Z \\ \hline
PS & PL & PS & Z & NS & NS \\ \hline
PL & PL & PS & Z & NS & NL \\ \hline
\end{tabular}
\end{table}

These rules ensure that the FLC dynamically adjusts the duty cycle to track the MPP under varying conditions. For example:
\begin{itemize}
    \item If \(\Delta P\) is \textit{Negative Large (NL)} and \(\Delta V\) is \textit{Negative Large (NL)}, then the output is \textit{Negative Large (NL)}.
    \item If \(\Delta P\) is \textit{Positive Small (PS)} and \(\Delta V\) is \textit{Zero (Z)}, then the output is \textit{Positive Small (PS)}.
    \item If \(\Delta P\) is \textit{Zero (Z)} and \(\Delta V\) is \textit{Zero (Z)}, then the output is \textit{Zero (Z)}.
\end{itemize}

\section{Simulation Results}
The performance of the proposed Fuzzy Logic-based Maximum Power Point Tracking (MPPT) system is rigorously evaluated through comprehensive simulations under dynamically changing solar irradiation conditions. The system is specifically designed to maintain stable operation and accurately track the maximum power point (MPP) even in the presence of significant variations in irradiation levels. This section presents detailed transient simulation results, highlighting the system's robustness and adaptability to fluctuating environmental conditions.

To assess the system's dynamic response, transient simulations are conducted under varying solar irradiation conditions. The irradiation level is systematically varied from 1000 W/m\(^2\) to 500 W/m\(^2\) and then back to 1000 W/m\(^2\), simulating real-world scenarios such as partial shading, cloud cover, or changes in weather. The results illustrate that the system effectively adjusts the duty cycle in real-time to maintain the MPP, ensuring optimal power transfer from the PV panel to the battery. The Fuzzy Logic Controller (FLC) demonstrates its ability to handle nonlinearities and uncertainties, providing a stable and efficient response to dynamic changes in irradiation.

\begin{figure}[ht]
\centering
\includegraphics[width=0.8\columnwidth]{tran.png}
\caption{Transient simulation results showing the system's response to dynamically changing solar irradiation. The duty cycle is dynamically adjusted to maintain the MPP, ensuring stable and efficient operation under varying conditions.}
\label{fig:transient_simulation}
\end{figure}

The simulation result (Figure 8) underscore the effectiveness and robustness of the proposed Fuzzy Logic-based MPPT system. The system not only tracks the MPP with high accuracy but also adapts seamlessly to sudden changes in solar irradiation. This adaptability is critical for real-world applications, where environmental conditions can vary rapidly and unpredictably. The ability to maintain stable operation under such conditions highlights the practical viability of the proposed system for use in photovoltaic battery charging applications.

\section{Conclusion}
This paper presents a Fuzzy Logic-based MPPT system for PV battery chargers using an advanced buck converter. The proposed system demonstrates improved efficiency and dynamic response compared to conventional MPPT methods. The buck converter is designed to operate at its maximum efficiency point, ensuring optimal power transfer from the PV panel to the battery. The Fuzzy Logic Controller dynamically adjusts the duty cycle to track the MPP under varying environmental conditions. Simulation results validate the effectiveness of the proposed approach, highlighting its potential for use in real-world PV battery charging systems.

\section*{Acknowledgements}
This work was supported by Your University and Your Funding Agency.

\end{document}